% ============================================================
% CHƯƠNG 5: KẾT LUẬN VÀ HƯỚNG PHÁT TRIỂN
% ============================================================

\chapter{Kết luận và hướng phát triển}
\label{chap:conclusion}

\section{Kết luận}
\label{sec:conclusion}

Khóa luận này đã giải quyết bài toán dự đoán mức độ stress liên tục (thang 1--9) từ dữ liệu cảm biến và hành vi người dùng thông qua một hệ thống học sâu đa mô-đun có nhận thức ngữ cảnh. Trên cơ sở tổng hợp lý thuyết về stress sinh lý học \cite{chrousos2009stress, mcewen2008central}, nhận dạng hoạt động từ cảm biến \cite{kwapisz2010activity, ordonez2016deep}, và mô hình chuỗi thời gian hồi quy (LSTM/Bi-LSTM) \cite{hochreiter1997lstm, schuster1997bidirectional}, khóa luận đã xây dựng và kiểm chứng thực nghiệm một pipeline hoàn chỉnh từ dữ liệu thô đến mô hình dự đoán.

Các kết quả chính đạt được:

\begin{enumerate}
    \item \textbf{Xây dựng thành công hệ thống end-to-end:}
    
    Hệ thống gồm 5 mô-đun liên kết chặt chẽ: (i) HAR, (ii) sinh dữ liệu đa phương thức có context-aware stress modifier, (iii) lựa chọn đặc trưng, (iv) pipeline chống rò rỉ dữ liệu, và (v) huấn luyện/tối ưu/đánh giá mô hình. Cấu trúc này cho phép giải quyết nhất quán toàn bộ bài toán thay vì xử lý rời rạc từng phần.

    \item \textbf{Mô-đun HAR đạt độ chính xác cao và ổn định:}
    
    Mô hình Stacked Bi-LSTM cho HAR đạt accuracy tổng thể 96.19\% trên tập WISDM, với F1-score cao ở các hoạt động chính (Jogging, Walking, Sitting, Standing). Kết quả này đảm bảo nhãn hoạt động đầu vào cho mô-đun stress prediction có độ tin cậy tốt.

    \item \textbf{Mô hình dự đoán stress đạt hiệu suất cao:}
    
    Trên tập 13 đặc trưng, mô hình Stacked Bi-LSTM baseline đạt R² = 0.9245, MAE = 0.6855. Sau tối ưu siêu tham số bằng Bayesian Optimization \cite{snoek2012practical}, mô hình tuned đạt R² = 0.9555 và MAE = 0.4414, tương ứng cải thiện 22.8\% MAE và 14.2\% RMSE so với baseline.

    \item \textbf{Tối ưu siêu tham số chứng minh hiệu quả rõ rệt:}
    
    Cùng một kiến trúc Bi-LSTM, chỉ bằng tối ưu hyperparameters đã cải thiện mạnh hiệu năng đồng thời giảm gần một nửa số tham số (320K xuống 164K). Kết quả này xác nhận rằng tuning là yếu tố quyết định, thậm chí quan trọng hơn việc thay đổi sang kiến trúc mới trong nhiều trường hợp.

    \item \textbf{Phân tích đặc trưng và lỗi cho thấy tính nhất quán khoa học:}
    
    Phân tích tầm quan trọng đặc trưng bằng 4 phương pháp (Permutation, SHAP, Correlation, RF Surrogate) cho thấy Mood\_Score và Heart\_Rate là hai đặc trưng quan trọng nhất, phù hợp với các nghiên cứu stress dựa trên tín hiệu sinh lý và hành vi \cite{hovsepian2015cstress, can2019stress}. Error analysis chỉ ra mô hình tuned cải thiện đặc biệt mạnh ở nhóm stress rất cao (8--9), giảm MAE 46.1\%, đây là nhóm có ý nghĩa ứng dụng thực tiễn cao trong cảnh báo sớm.

    \item \textbf{So sánh mô hình cho kết luận thuyết phục:}
    
    Trên cùng pipeline dữ liệu và cùng điều kiện huấn luyện, mô hình Stacked Bi-LSTM tuned vượt trội so với MLP, Simple LSTM, Stacked Bi-LSTM baseline và Stacked Bi-GRU về cả độ chính xác (MAE/RMSE/R²) lẫn hiệu quả tham số. Kết quả này khẳng định lựa chọn kiến trúc đề xuất là phù hợp.
\end{enumerate}

Tổng hợp các kết quả trên cho thấy mục tiêu nghiên cứu đã được hoàn thành: khóa luận không chỉ chứng minh tính khả thi của việc dự đoán stress liên tục bằng Deep Learning, mà còn đề xuất được một quy trình có thể tái sử dụng cho các bài toán sức khỏe số dựa trên dữ liệu cảm biến và hành vi.

\section{Đóng góp học thuật và điểm khác biệt}
\label{sec:academic_contributions}

Để nhấn mạnh bản chất nghiên cứu thay vì báo cáo triển khai kỹ thuật, các đóng góp của khóa luận có thể được nhìn nhận theo ba lớp học thuật sau:

\begin{enumerate}
    \item \textbf{Đóng góp ở mức bài toán (problem formulation):}
    
    Khóa luận chuyển trọng tâm từ phân loại stress rời rạc (binary/multi-class) sang \textbf{hồi quy stress liên tục} trên thang 1--9. Đây là thay đổi quan trọng vì cho phép biểu diễn cường độ stress mịn hơn, phù hợp với mục tiêu theo dõi và can thiệp sớm.

    \item \textbf{Đóng góp ở mức phương pháp (methodological contribution):}
    
    Khóa luận đề xuất khung \textit{context-aware multi-modal} tích hợp HAR với mô hình chuỗi thời gian Bi-LSTM và cơ chế Context-Stress Modifiers để giảm nhập nhằng sinh lý. Điểm khác biệt so với nhiều công trình cũ là không xử lý tín hiệu sinh lý như biến độc lập, mà diễn giải theo tổ hợp hoạt động-ngữ cảnh-thời gian.

    \item \textbf{Đóng góp ở mức bằng chứng thực nghiệm (empirical evidence):}
    
    Nghiên cứu không dừng ở một mô hình đơn lẻ mà thực hiện so sánh công bằng 5 kiến trúc, kết hợp 4 phương pháp giải thích đặc trưng và phân tích lỗi đa chiều. Kết quả cho thấy Bi-LSTM tuned vượt trội đồng thời về độ chính xác và hiệu quả tham số, từ đó tạo bằng chứng thuyết phục cho lựa chọn phương pháp.
\end{enumerate}

Từ góc nhìn học thuật, giá trị cốt lõi của khóa luận không chỉ nằm ở con số R² cao, mà ở việc xác lập được mối liên hệ nhất quán giữa \textbf{câu hỏi nghiên cứu $\rightarrow$ thiết kế phương pháp $\rightarrow$ bằng chứng thực nghiệm $\rightarrow$ kết luận khoa học}.

\section{Hạn chế của khóa luận}
\label{sec:limitations}

Mặc dù đạt kết quả tích cực, khóa luận vẫn tồn tại một số hạn chế. Mục này đồng thời tổng hợp các đe dọa đến tính hiệu lực nghiên cứu dưới dạng các giới hạn học thuật và triển khai:

\begin{enumerate}
    \item \textbf{Hạn chế về construct validity (độ giá trị cấu trúc):}
    
    Stress\_Level được sinh theo mô hình mô phỏng có căn cứ (semi-synthetic), chưa phải nhãn lâm sàng đo trực tiếp từ nghiên cứu thực địa đa đối tượng. Vì vậy, kết quả phản ánh tốt trong khung giả định nhưng chưa đại diện đầy đủ cho thực tế lâm sàng.

    \item \textbf{Hạn chế về internal validity (độ giá trị nội tại):}
    
    Pipeline đã kiểm soát data leakage theo thứ tự Split $\to$ Encode $\to$ Normalize $\to$ Sequence, tuy nhiên vẫn tồn tại rủi ro thiên lệch do các giả định của bộ sinh dữ liệu và chưa có đầy đủ thí nghiệm ablation cho từng thành phần ngữ cảnh.

    \item \textbf{Hạn chế về external validity và đa dạng đối tượng:}
    
    Kịch bản mô phỏng tập trung vào hồ sơ người dùng điển hình trong 30 ngày; mô hình chưa được kiểm chứng trên tập dữ liệu độc lập bên ngoài (external validation set) hoặc subject-wise cross-validation. Do đó, khả năng tổng quát hóa theo giới, tuổi, nghề nghiệp, bệnh nền và lối sống vẫn cần xác nhận thêm.

    \item \textbf{Hạn chế về conclusion validity (độ giá trị kết luận):}
    
    Kết luận hiện chủ yếu dựa trên các chỉ số trung bình (MAE/RMSE/R²). Nghiên cứu chưa thực hiện đầy đủ kiểm định thống kê suy luận (ví dụ paired test, bootstrap confidence interval) qua nhiều lần chạy độc lập để tăng độ chắc chắn cho so sánh mô hình.

    \item \textbf{Hạn chế về không gian mô hình:}
    
    Phần so sánh tập trung vào nhóm LSTM family (LSTM, Bi-LSTM, Bi-GRU) và MLP, chưa mở rộng sang Transformer/Attention-based time-series models, Temporal Convolutional Networks hoặc kiến trúc lai đa nhánh.

    \item \textbf{Hạn chế về nguồn lực tính toán:}
    
    Toàn bộ thí nghiệm được huấn luyện trên CPU, dẫn đến giới hạn số lượng trial trong hyperparameter tuning và giới hạn quy mô thử nghiệm (số mô hình, số cấu hình, số lần lặp).

    \item \textbf{Hạn chế về triển khai thời gian thực:}
    
    Khóa luận dừng ở mức nghiên cứu mô hình ngoại tuyến (offline). Các yêu cầu triển khai thực tế như tối ưu latency, tiêu thụ pin, drift detection và cập nhật mô hình online chưa được xử lý đầy đủ.
\end{enumerate}

Những hạn chế trên không phủ nhận giá trị khoa học của kết quả đạt được, nhưng chỉ ra ranh giới áp dụng hiện tại và là cơ sở để xây dựng lộ trình phát triển có kiểm chứng mạnh hơn trong tương lai.

\section{Hướng phát triển}
\label{sec:future_work}

Từ kết quả và hạn chế đã phân tích, các hướng phát triển được đề xuất theo ba mức thời gian:

\subsection{Hướng ngắn hạn (3--6 tháng)}

\begin{enumerate}
    \item \textbf{Thu thập dữ liệu thực tế trên smartphone/wearable:}
    
    Xây dựng ứng dụng thu thập dữ liệu theo consent để ghi nhận đồng thời accelerometer, heart rate, screen usage, location và self-report stress theo thời gian. Mục tiêu là tạo tập dữ liệu thực địa đầu tiên để kiểm chứng mô hình.

    \item \textbf{Chuẩn hóa giao thức gán nhãn stress:}
    
    Kết hợp self-report ngắn (EMA -- Ecological Momentary Assessment) với các chỉ dấu sinh lý để giảm sai lệch nhãn, tăng độ tin cậy của ground-truth.

    \item \textbf{Mở rộng benchmarking mô hình:}
    
    Bổ sung Transformer và Attention-based architectures để so sánh trực tiếp với Bi-LSTM tuned trên cùng pipeline dữ liệu.

    \item \textbf{Tối ưu bộ đặc trưng trong điều kiện thực tế:}
    
    Đánh giá lại 13 đặc trưng trên dữ liệu thực để xem đặc trưng nào bền vững theo cá thể, đặc trưng nào cần hiệu chỉnh hoặc loại bỏ.
\end{enumerate}

\subsection{Hướng trung hạn (6--18 tháng)}

\begin{enumerate}
    \item \textbf{Xây dựng hệ thống dự đoán stress thời gian thực:}
    
    Triển khai pipeline suy luận online trên thiết bị di động, tối ưu tốc độ và năng lượng, cung cấp cảnh báo stress theo ngữ cảnh gần thời gian thực.

    \item \textbf{Cá nhân hóa mô hình (personalization):}
    
    Áp dụng transfer learning hoặc meta-learning để tinh chỉnh từ mô hình nền sang mô hình cá nhân với số lượng mẫu nhỏ cho từng người dùng.

    \item \textbf{Học liên kết bảo mật (Federated Learning):}
    
    Huấn luyện mô hình phân tán trên thiết bị người dùng để tăng riêng tư dữ liệu cá nhân, giảm nhu cầu tập trung dữ liệu nhạy cảm.

    \item \textbf{Bổ sung cơ chế giải thích cho người dùng cuối:}
    
    Tích hợp giải thích theo phiên (session-level explanation) dựa trên SHAP để giúp người dùng hiểu yếu tố nào đang làm tăng stress ở từng thời điểm.
\end{enumerate}

\subsection{Hướng dài hạn (trên 18 tháng)}

\begin{enumerate}
    \item \textbf{Mở rộng dữ liệu đa cảm biến sinh học:}
    
    Tích hợp thêm EDA (electrodermal activity), HRV, EEG hoặc marker hormone gián tiếp để tăng độ tin cậy của dự đoán stress đa chiều.

    \item \textbf{Từ dự đoán sang can thiệp (prediction-to-intervention):}
    
    Xây dựng hệ thống đề xuất can thiệp cá nhân hóa (nhắc nghỉ ngắn, bài thở, đi bộ nhẹ, giảm screen load) dựa trên mức stress dự đoán và lịch sử phản hồi người dùng.

    \item \textbf{Mô hình sức khỏe tâm thần tổng quát:}
    
    Mở rộng từ dự đoán stress đơn biến sang mô hình đa mục tiêu: stress, chất lượng giấc ngủ, nguy cơ kiệt sức (burnout), trạng thái cảm xúc và mức tập trung.

    \item \textbf{Đánh giá lâm sàng và ứng dụng thực tiễn:}
    
    Hợp tác với chuyên gia tâm lý/y khoa để triển khai nghiên cứu theo dõi dọc (longitudinal study), đánh giá tác động của hệ thống trong bối cảnh chăm sóc sức khỏe cộng đồng.
\end{enumerate}

\section{Tổng kết chương}
\label{sec:chapter5_summary}

Chương 5 đã tổng kết toàn bộ kết quả nghiên cứu, xác nhận rằng mục tiêu chính của khóa luận đã đạt được: xây dựng thành công hệ thống dự đoán stress liên tục dựa trên dữ liệu cảm biến và hành vi với hiệu suất cao (R² = 0.9555, MAE = 0.4414). Đồng thời, chương này cũng chỉ ra các giới hạn hiện tại và đề xuất lộ trình phát triển theo từng giai đoạn, tạo nền tảng để chuyển từ nghiên cứu mô hình ngoại tuyến sang ứng dụng sức khỏe số thời gian thực trong thực tế.